\sectionguard
\section*{The Appointed Formulary in Full}

Everything the 1962 \latin{Missale Romanum} prints for this Mass, in the order
it prints it: the Latin of the Vatican \latin{editio typica} exactly as the
controlling facsimile has it, each element under its own printed heading,
scriptural reference and marginal number, and beneath it the English of an
identified public-domain witness. This is the canonical full research edition;
a reader who came for the texts has them here.

\textbf{Formulary.} \latin{DOMINICA DUODECIMA post Pentecosten}, \latin{II
classis}, printed pp.~392--393, marginal nos.\ 1552--1561. Green follows the
general rubrics for Sundays after Pentecost and is not printed on those pages.
The formulary begins immediately after the Postcommunion of the Eleventh
Sunday (no.\ 1551) and ends immediately before the heading \latin{DOMINICA
DECIMA TERTIA post Pentecosten} (no.\ 1562). Nothing between those boundaries
is omitted here. It appoints no Tract, Sequence, second oration, blessing or
ritual text, and no second Collect, Secret or Postcommunion.

\textbf{The Latin.} Received liturgical text of the Roman Rite, older than the
edition that prints it. No public-domain conclusion is claimed for the 1962
typical edition itself; what permits the Latin here is 17 U.S.C.\ \S103(b),
under which a revised edition's copyright reaches only that edition's own
contribution and not the preexisting material it reprints. All ten appointed
elements of this formulary are printed in the Pustet Missal of Ratisbon, 1862,
and located in the Venice Missal of 1570, both of which this repository tracks
as public-domain text; the terminal appendix states the basis and its bounds
in full. The project claims no rights in this text and has not composed,
translated, modernised, conflated or adapted any of it. Every form below was
visually collated at 300 and 600 dpi against the CMAA facsimile of the Vatican
typical edition, and independently against page images of the 1962 Benziger
\latin{editio iuxta typicam}. Where the two editions differ, the Vatican
typical reading is what appears here and the difference is tabulated in the
detailed commentary.

\textbf{The English.} Scriptural elements are the Douay--Rheims as revised by
Bishop Challoner (1749--1752) --- the English of the Clementine Vulgate the
missal prints --- in the Project Gutenberg e-text of that revision. The three
orations are the anonymous English of \work{The Roman Missal translated into
the English language for the use of the laity}, first revised edition
(Philadelphia: Eugene Cummiskey, 1861), from the three lines that book
supplies for this formulary, reproduced with its own nineteenth-century
spelling and its abbreviated conclusion ``Thro'.'' Both are public domain in
the United States. Neither is an official liturgical translation, and neither
is the text of the 1962 books: this is a study companion, not an altar book or
a hand missal. The project has translated nothing. Two chants of this
formulary follow the old chant psalter or the Old Latin rather than the
Clementine Vulgate, so the registered English does not answer their exact
wording at several points; each gap is marked at the element and left open.

\elementguard
\subsection*{1. Introit \cue{Int.}}

\begin{appointedlatin}{\latin{Ant. ad Introitum} --- Ps. 69, 2-3 --- marginal no.\ 1552}
Deus, in adiutórium meum inténde: Dómine, ad adiuvándum me festína:
confundántur et revereántur inimíci mei, qui quærunt ánimam meam.
\textnormal{Ps. ibid., 4} Avertántur retrórsum, et erubéscant: qui cógitant
mihi mala. ℣.~Glória Patri.
\end{appointedlatin}

\begin{namedtranslation}{English: Douay--Rheims (Challoner), Ps. 69:2--3 and 69:4}
O God, come to my assistance; O Lord, make haste to help me. Let them be
confounded and ashamed that seek my soul. \textit{Ps.} Let them be turned
backward, and blush for shame that desire evils to me.
℣.~\latin{Glória Patri.}
\end{namedtranslation}

\englishgap{The antiphon and verse follow the chant psalter, not the
Clementine the Douay translates, and the English above therefore misses the
missal at two points. The antiphon sings \latin{confundántur et revereántur
inimíci mei}, ``my enemies''; Clementine Ps.\ 69:3 has no \latin{inimici mei},
so the Douay's ``that seek my soul'' hangs from no named enemy. The verse
sings \latin{qui cógitant mihi mala}, ``who devise evils against me''; the
Clementine, and therefore the Douay, reads \latin{qui volunt mihi mala},
``that desire evils to me'' --- desiring is what the English answers, devising
is what the missal prints. The verse also stops at the half-verse; the
Clementine's continuation (\latin{avertantur statim erubescentes\dots}) is not
appointed. The doxology cue \latin{Glória Patri} belongs to the Ordinary
rather than to this formulary and is a scriptureless liturgical formula; no
English is supplied for it here.}

\elementguard
\subsection*{2. Collect \cue{Coll.}}

\begin{appointedlatin}{\latin{Oratio} --- marginal no.\ 1553}
Omnípotens et miséricors Deus, de cuius múnere venit, ut tibi a fidélibus tuis
digne et laudabíliter serviátur: tríbue, quǽsumus, nobis; ut ad promissiónes
tuas sine offensióne currámus. Per Dóminum.
\end{appointedlatin}

\begin{namedtranslation}{English: Cummiskey 1861, Collect of this formulary, printed p.\ 420}
O Almighty and merciful God, from whose gift it proceedeth that thy people
worthily serve thee; grant we beseech thee, that we may run on, without
stumbling, to the obtaining the effects of thy promises. Thro'.
\end{namedtranslation}

\englishgap{The 1861 rendering gives one adverb where the Latin prays two:
``worthily serve thee'' answers \latin{digne} and leaves \latin{et
laudabíliter} --- ``and laudably,'' service that is also praise --- without a
separate English word. It also renders \latin{a fidélibus tuis}, ``by thy
faithful,'' as ``thy people.'' The grammatical argument in the commentary is
therefore conducted on the Latin.}

\elementguard
\subsection*{3. Epistle \cue{Ep.}}

\begin{appointedlatin}{\latin{Léctio Epístolæ beáti Pauli Apóstoli ad Corínthios.} --- 2 Cor. 3, 4-9 --- marginal no.\ 1554}
Fratres: Fidúciam talem habémus per Christum ad Deum: non quod sufficiéntes
simus cogitáre áliquid a nobis, quasi ex nobis: sed sufficiéntia nostra ex Deo
est: qui et idóneos nos fecit minístros novi testaménti: non líttera, sed
spíritu: líttera enim occídit, spíritus autem vivíficat. Quod si ministrátio
mortis, lítteris deformáta in lapídibus, fuit in glória; ita ut non possent
inténdere fílii Israël in fáciem Móysi, propter glóriam vultus eius, quæ
evacuátur: quómodo non magis ministrátio Spíritus erit in glória? Nam si
ministrátio damnatiónis glória est: multo magis abúndat ministérium iustítiæ
in glória.
\end{appointedlatin}

\begin{namedtranslation}{English: Douay--Rheims (Challoner), 2 Cor. 3:4--9}
And such confidence we have, through Christ, towards God. Not that we are
sufficient to think any thing of ourselves, as of ourselves: but our
sufficiency is from God. Who also hath made us fit ministers of the new
testament, not in the letter but in the spirit. For the letter killeth: but
the spirit quickeneth. Now if the ministration of death, engraven with letters
upon stones, was glorious (so that the children of Israel could not
steadfastly behold the face of Moses, for the glory of his countenance), which
is made void: How shall not the ministration of the spirit be rather in glory?
For if the ministration of condemnation be glory, much more the ministration
of justice aboundeth in glory.
\end{namedtranslation}

\englishgap{The liturgical incipit \latin{Fratres:} is the missal's own and is
not translated; it displaces the connective of Clementine 2 Cor.\ 3:4,
\latin{Fidúciam \textbf{autem} talem habémus} --- the Douay's opening ``And''
renders the \latin{autem} that the lesson drops. The pericope ends at v.~9;
v.~10 is outside the appointed lesson.}

\elementguard
\subsection*{4. Gradual \cue{Grad.}}

\begin{appointedlatin}{\latin{Graduale} --- Ps. 33, 2-3 --- marginal no.\ 1555}
Benedícam Dóminum in omni témpore: semper laus eius in ore meo. ℣.~In Dómino
laudábitur ánima mea: áudiant mansuéti, et læténtur.
\end{appointedlatin}

\begin{namedtranslation}{English: Douay--Rheims (Challoner), Ps. 33:2--3}
I will bless the Lord at all times, his praise shall be always in my mouth.
℣.~In the Lord shall my soul be praised: let the meek hear and rejoice.
\end{namedtranslation}

\elementguard
\subsection*{5. Alleluia \cue{All.}}

\begin{appointedlatin}{Allelúia --- Ps. 87, 2 --- marginal no.\ 1556}
Allelúia, allelúia. ℣.~Dómine, Deus salútis meæ, in die clamávi et nocte coram
te. Allelúia.
\end{appointedlatin}

\begin{namedtranslation}{English: Douay--Rheims (Challoner), Ps. 87:2}
O Lord, the God of my salvation: I have cried in the day, and in the night
before thee. \latin{Allelúia.}
\end{namedtranslation}

\elementguard
\subsection*{6. Gospel \cue{Gosp.}}

\begin{appointedlatin}{✠ \latin{Sequéntia sancti Evangélii secúndum Lucam.} --- Luc. 10, 23-37 --- marginal no.\ 1557}
In illo témpore: Dixit Iesus discípulis suis: Beáti óculi, qui vident quæ vos
vidétis. Dico enim vobis, quod multi prophétæ et reges voluérunt vidére quæ
vos vidétis, et non vidérunt: et audíre quæ audítis, et non audiérunt. Et ecce
quidam legisperítus surréxit, tentans illum, et dicens: Magíster, quid
faciéndo vitam ætérnam possidébo? At ille dixit ad eum: In lege quid scriptum
est? quómodo legis? Ille respóndens, dixit: Díliges Dóminum Deum tuum ex toto
corde tuo, et ex tota ánima tua, et ex ómnibus víribus tuis, et ex omni mente
tua: et próximum tuum sicut teípsum. Dixítque illi: Recte respondísti: hoc
fac, et vives. Ille autem volens iustificáre seípsum, dixit ad Iesum: Et quis
est meus próximus? Suscípiens autem Iesus, dixit: Homo quidam descendébat ab
Ierúsalem in Iéricho, et íncidit in latrónes, qui étiam despoliavérunt eum: et
plagis impósitis abiérunt, semivívo relícto. Accidit autem ut sacérdos quidam
descénderet eádem via: et viso illo præterívit. Simíliter et levíta, cum esset
secus locum, et vidéret eum, pertránsiit. Samaritánus autem quidam iter
fáciens, venit secus eum: et videns eum, misericórdia motus est. Et apprópians
alligávit vúlnera eius, infúndens óleum et vinum: et impónens illum in
iuméntum suum, duxit in stábulum, et curam eius egit. Et áltera die prótulit
duos denários, et dedit stabulário, et ait: Curam illíus habe: et quodcúmque
supererogáveris, ego cum redíero, reddam tibi. Quis horum trium vidétur tibi
próximus fuísse illi, qui íncidit in latrónes? At ille dixit: Qui fecit
misericórdiam in illum. Et ait illi Iesus: Vade, et tu fac simíliter.
\end{appointedlatin}

\appointedrubric{Credo.}{printed at the right-hand end of the Gospel's final
line}

\begin{namedtranslation}{English: Douay--Rheims (Challoner), Luke 10:23--37}
Blessed are the eyes that see the things which you see. For I say to you that
many prophets and kings have desired to see the things that you see and have
not seen them; and to hear the things that you hear and have not heard them.
And behold a certain lawyer stood up, tempting him and saying, Master, what
must I do to possess eternal life? But he said to him: What is written in the
law? How readest thou? He answering, said: Thou shalt love the Lord thy God
with thy whole heart and with thy whole soul and with all thy strength and
with all thy mind: and thy neighbour as thyself. And he said to him: Thou hast
answered right. This do: and thou shalt live. But he willing to justify
himself, said to Jesus: And who is my neighbour? And Jesus answering, said: A
certain man went down from Jerusalem to Jericho and fell among robbers, who
also stripped him and having wounded him went away, leaving him half dead. And
it chanced, that a certain priest went down the same way: and seeing him,
passed by. In like manner also a Levite, when he was near the place and saw
him, passed by. But a certain Samaritan, being on his journey, came near him:
and seeing him, was moved with compassion: And going up to him, bound up his
wounds, pouring in oil and wine: and setting him upon his own beast, brought
him to an inn and took care of him. And the next day he took out two pence and
gave to the host and said: Take care of him; and whatsoever thou shalt spend
over and above, I, at my return, will repay thee. Which of these three, in thy
opinion, was neighbour to him that fell among the robbers? But he said: He
that shewed mercy to him. And Jesus said to him: Go, and do thou in like
manner.
\end{namedtranslation}

\englishgap{The liturgical incipit \latin{In illo témpore: Dixit Iesus
discípulis suis:} is the missal's own and is not translated. It displaces
Luke's narrative frame, Clementine 10:23 \latin{Et convérsus ad discípulos
suos, dixit:} --- the evangelist's ``And turning to his disciples, he said''
is outside the appointed text, so the Douay's opening words for v.~23 are not
printed above.}

\elementguard
\subsection*{7. Offertory \cue{Off.}}

\begin{appointedlatin}{\latin{Antiphona ad Offertorium} --- Exodi 32, 11, 13 et 14 --- marginal no.\ 1558}
Precátus est Móyses in conspéctu Dómini Dei sui, et dixit: Quare, Dómine,
irásceris in pópulo tuo? Parce iræ ánimæ tuæ: meménto Abraham, Isaac, et
Iacob, quibus iurásti dare terram fluéntem lac et mel. Et placátus factus est
Dóminus de malignitáte, quam dixit fácere pópulo suo.
\end{appointedlatin}

\begin{namedtranslation}{English: Douay--Rheims (Challoner), Ex. 32:11, 13 and 14, printed whole}
\textit{v.~11.} But Moses besought the Lord his God, saying: Why, O Lord, is
thy indignation enkindled against thy people, whom thou hast brought out of
the land of Egypt, with great power, and with a mighty hand?

\textit{v.~13.} Remember Abraham, Isaac, and Israel, thy servants, to whom
thou sworest by thy own self, saying: I will multiply your seed as the stars
of heaven: and this whole land that I have spoken of, I will give to your
seed, and you shall possess it for ever:

\textit{v.~14.} And the Lord was appeased from doing the evil which he had
spoken against his people.
\end{namedtranslation}

\englishgap{This chant is an Old-Latin adaptation of the narrative, not the
Vulgate the Douay translates, so the three verses are printed whole above and
no English sentence is stitched to match the chant. Four gaps stand. The
chant's \latin{Parce iræ ánimæ tuæ} carries the matter of v.~12 --- a verse
the missal's own citation does not name --- and has no counterpart in any
cited verse; the Douay's v.~12, for the record, reads ``let thy anger cease,
and be appeased upon the wickedness of thy people.'' The chant sings
\latin{Iacob} where v.~13, and therefore the Douay, reads ``Israel, thy
servants.'' The chant's oath is \latin{dare terram fluéntem lac et mel}, ``to
give a land flowing with milk and honey,'' where v.~13's oath is the
stars-and-seed promise printed above --- the milk-and-honey formula belongs to
Ex.\ 3:8 and 33:3, not to this verse. And the chant's opening \latin{in
conspéctu Dómini Dei sui}, ``in the sight of the Lord his God,'' answers to
nothing in v.~11's ``besought the Lord his God.''}

\elementguard
\subsection*{8. Secret \cue{Sec.}}

\begin{appointedlatin}{\latin{Secreta} --- marginal no.\ 1559}
Hóstias, quǽsumus, Dómine, propítius inténde, quas sacris altáribus
exhibémus: ut, nobis indulgéntiam largiéndo, tuo nómini dent honórem. Per
Dóminum nostrum.
\end{appointedlatin}

\appointedrubric{Præfatio de Ssma Trinitate.}{the Preface of the Most Holy
Trinity; \latin{Ssma} printed under a tilde abbreviation for
\latin{Sanctissima}}

\begin{namedtranslation}{English: Cummiskey 1861, Secret of this formulary, printed p.\ 422}
Mercifully look down, O Lord, on the offerings we lay on thy holy altar; that
they may be to the honour of thy name, by obtaining pardon for us. Thro'.
\end{namedtranslation}

\englishgap{The 1861 book prints no incipit for this Secret, so its identity
was confirmed clause against clause --- ``the offerings we lay on thy holy
altar'' against \latin{quas sacris altáribus exhibémus}, ``to the honour of
thy name, by obtaining pardon for us'' against \latin{nobis indulgéntiam
largiéndo, tuo nómini dent honórem} --- and by reading the orations of the
eleventh and thirteenth Sundays on either side to prove the filing does not
slip by one. The English reverses the Latin's order of honour and pardon and
makes ``look down'' of \latin{propítius inténde}, ``favourably attend'';
the commentary conducts the exchange-structure argument on the Latin.}

\elementguard
\subsection*{9. Communion \cue{Comm.}}

\begin{appointedlatin}{\latin{Antiphona ad Communionem} --- Ps. 103, 13 et 14-15 --- marginal no.\ 1560}
De fructu óperum tuórum, Dómine, satiábitur terra: ut edúcas panem de terra,
et vinum lætíficet cor hóminis: ut exhílaret fáciem in óleo, et panis cor
hóminis confírmet.
\end{appointedlatin}

\begin{namedtranslation}{English: Douay--Rheims (Challoner), Ps. 103:13--15, printed whole}
\textit{v.~13.} Thou waterest the hills from thy upper rooms: the earth shall
be filled with the fruit of thy works:

\textit{v.~14.} Bringing forth grass for cattle, and herb for the service of
men. That thou mayst bring bread out of the earth:

\textit{v.~15.} And that wine may cheer the heart of man. That he may make the
face cheerful with oil: and that bread may strengthen man's heart.
\end{namedtranslation}

\englishgap{The antiphon takes v.~13's second clause with its word order
inverted and a vocative added --- \latin{De fructu óperum tuórum,
\textbf{Dómine}, satiábitur terra} --- and the \latin{Dómine} is in no verse
of the psalm here, so the Douay has no word for it. The antiphon then skips
v.~14's first clause, exactly as the printed citation \latin{13 et 14-15}
declares, and runs continuously from \latin{ut edúcas} to the end of v.~15.
The three verses are printed whole above with the antiphon's path through them
stated; no English join is composed.}

\elementguard
\subsection*{10. Postcommunion \cue{Postcomm.}}

\begin{appointedlatin}{\latin{Postcommunio} --- marginal no.\ 1561}
Vivíficet nos, quǽsumus, Dómine, huius participátio sancta mystérii: et
páriter nobis expiatiónem tríbuat, et munímen. Per Dóminum.
\end{appointedlatin}

\begin{namedtranslation}{English: Cummiskey 1861, Postcommunion of this formulary, printed p.\ 423}
May the sacred participation of these thy mysteries, O Lord, we beseech thee,
give us life; and be to us both an expiation and protection. Thro'.
\end{namedtranslation}

\englishgap{The Latin asks one subject to \latin{tríbuat} --- to
\textit{bestow} --- two named gifts, \latin{expiatiónem \dots\ et munímen},
with \latin{páriter} binding them as equal and simultaneous; the 1861 English
flattens the bestowal into ``be to us both,'' and gives no word for
\latin{páriter}. It also pluralises: ``these thy mysteries'' where the Latin
has the singular \latin{huius \dots\ mystérii}. The commentary's argument
about the pair is conducted on the Latin.}

\clearpage
